% !TeX root = ../main.tex
\section{Cache-Miss Impact}\label{sec:cache-miss}

\begin{hypothesis}
  As the SAXPY footprint grows past each cache, time per element should rise with the misses per element.
  Sequential order should rise gently (the prefetcher turns misses into hits), random order sharply. Because
  independent iterations overlap their misses, the effective DRAM miss cost should be well below the
  zero-queue latency.
\end{hypothesis}

\begin{knobs}
  Kernel      & $y_i = a\,x_i + y_i$ (float64), two arrays, \qty{4}{\kibi\byte} pages \\
  Varied      & $N = 2^{8} \ldots 2^{24}$ (footprint $16N$ bytes sequential; $20N$ random, which adds a 4-byte index per element); order (sequential, random permutation) \\
  Counters    & L1D read misses, LLC read references/misses, dTLB walks, cycles, instructions (user space, per element) \\
  Script      & \code{Cache_and_Memory_Performance_Profiling/scripts/run.sh cache_miss} \\
\end{knobs}

\begin{figure}[tbp]
  \centering
  \begin{subfigure}[t]{0.49\linewidth}
  \IfDataExists{cache_miss_ctr}{%
    \begin{tikzpicture}
      \begin{axis}[acs, xmode=log, log basis x=2, ymode=log, width=0.97\linewidth, height=0.72\linewidth,
                   xlabel={$N$ (elements)}, ylabel={Time per element (ns)}, legend pos=north west,
                   legend style={font=\scriptsize}]
        \addplot+[iqr bars] table[x=n, y=ns_per_elem_median, iqr of=ns_per_elem, only if={pattern}{sequential}] {data/cache_miss_ctr.csv};
        \addlegendentry{sequential}
        \addplot+[iqr bars] table[x=n, y=ns_per_elem_median, iqr of=ns_per_elem, only if={pattern}{random}] {data/cache_miss_ctr.csv};
        \addlegendentry{random}
        \IfFileExists{data/amat_fit.csv}{%
        \addplot[only marks, mark=o, mark size=2.5pt, black] table[x=n, y=predicted_ns, col sep=comma] {data/amat_fit.csv};
        \addlegendentry{AMAT fit}}{}
        \acsvline{2048}{L1d}
        \acsvline{16384}{L2}
        \acsvline{1310720}{L3}
      \end{axis}
    \end{tikzpicture}}
    \caption{Runtime and fitted AMAT model.}
  \end{subfigure}\hfill
  \begin{subfigure}[t]{0.49\linewidth}
  \IfDataExists{cache_miss_ctr}{%
    \begin{tikzpicture}
      \begin{axis}[acs, xmode=log, log basis x=2, height=0.72\linewidth, ymin=0,
                   xlabel={$N$ (elements)}, ylabel={Events per element}, legend pos=north west,
                   legend style={font=\scriptsize}]
        \addplot+[mark=*, iqr bars] table[x=n, y=ctr_l1d_miss_pe_median, iqr of=ctr_l1d_miss_pe, only if={pattern}{random}] {data/cache_miss_ctr.csv};
        \addlegendentry{L1D miss (rand)}
        \addplot+[mark=triangle*, iqr bars] table[x=n, y=ctr_llc_miss_pe_median, iqr of=ctr_llc_miss_pe, only if={pattern}{random}] {data/cache_miss_ctr.csv};
        \addlegendentry{LLC miss (rand)}
        \addplot+[mark=diamond*, iqr bars] table[x=n, y=ctr_dtlb_miss_pe_median, iqr of=ctr_dtlb_miss_pe, only if={pattern}{random}] {data/cache_miss_ctr.csv};
        \addlegendentry{dTLB walk (rand)}
        \addplot+[mark=o, dashed, iqr bars] table[x=n, y=ctr_l1d_miss_pe_median, iqr of=ctr_l1d_miss_pe, only if={pattern}{sequential}] {data/cache_miss_ctr.csv};
        \addlegendentry{L1D miss (seq)}
        \addplot+[mark=x, dashed, iqr bars] table[x=n, y=ctr_llc_miss_pe_median, iqr of=ctr_llc_miss_pe, only if={pattern}{sequential}] {data/cache_miss_ctr.csv};
        \addlegendentry{LLC miss (seq)}
        \acsvline{2048}{L1d}
        \acsvline{16384}{L2}
        \acsvline{1310720}{L3}
      \end{axis}
    \end{tikzpicture}}
    \caption{Counters per element.}
  \end{subfigure}
  \caption{SAXPY time per element and counters vs.\ footprint and traversal order (two \code{double}
  arrays, $16N$ bytes; lines mark where $16N$ crosses L1/L2/L3; random order also reads a $4N$-byte index
  array, so its $20N$-byte footprint crosses each capacity at $0.8\times$ the marked $N$). Open circles: AMAT model fitted to the
  counters (\cref{eq:amat-fit}).}\label{fig:cache-miss}
\end{figure}

\subsection{Analysis}
Sequential SAXPY rises from \valqty{cm_seq_min_ns}{\nano\second} per element in L1 to
\valqty{cm_seq_max_ns}{\nano\second} in DRAM, a \val{cm_seq_span}$\times$ span, with only
\val{cm_seq_max_llc_miss} demand LLC misses per element at the largest size. Random traversal rises from \valqty{cm_rand_min_ns}{\nano\second} to
\valqty{cm_rand_max_ns}{\nano\second}, a \val{cm_rand_span}$\times$ span, with \val{cm_rand_max_l1d_miss}
L1D misses, \val{cm_rand_max_llc_miss} LLC misses and \val{cm_rand_max_dtlb_miss} page walks per element at
$N=2^{24}$ (a line of $x$, a line of $y$ and page-walk reads). At DRAM sizes the random curve is
\val{cm_dram_gap}$\times$ the sequential one: that gap is the miss cost the prefetcher hides.

\subsection{AMAT fitted to measured miss counts}
Classic AMAT, $t_{L1} + m_{L1}\bigl(t_{L2} + m_{L2}(t_{L3} + m_{L3}\,t_{\text{DRAM}})\bigr)$, assumes
serialised misses. We instead fit its per-element form to the counters, separately for each order:
\begin{equation}\label{eq:amat-fit}
  t_{\text{elem}} = t_0 + a_2\,h_{L2} + a_3\,h_{L3} + a_D\,m_D ,
\end{equation}
with per-element hits served by L2 ($h_{L2}$ = L1D misses $-$ LLC references), by L3 ($h_{L3}$ = LLC
references $-$ LLC misses) and DRAM misses ($m_D$ = LLC misses); $t_0$ absorbs L1-hit work. Least squares
over the 17 footprints gives, for random order,
$t_0=\valqty{amat_rand_t0}{\nano\second}$, $a_2=\valqty{amat_rand_a2}{\nano\second}$,
$a_3=\valqty{amat_rand_a3}{\nano\second}$, $a_D=\valqty{amat_rand_aD}{\nano\second}$ ($R^2=\val{amat_rand_r2}$);
for sequential order $a_D=\valqty{amat_seq_aD}{\nano\second}$ ($R^2=\val{amat_seq_r2}$). The fitted $a_D$ is
the \emph{exposed} cost of one DRAM miss; the zero-queue DRAM latency divided by it gives the average
number of overlapped misses: \val{amat_rand_mlp} for random and \val{amat_seq_mlp} for sequential SAXPY
(whose remaining misses come in long, predictable runs). For random order the zero-queue L2 and L3 latencies
are likewise \val{amat_rand_ovlL2}$\times$ and \val{amat_rand_ovlL3}$\times$ the fitted $a_2$ and $a_3$. All
three factors are close to the \valnum{mlp_t1_Lmem_max} demand misses one core keeps in flight
(\cref{sec:intensity}). Classic serialised AMAT confirms this at $N=2^{24}$: with the
\cref{tab:zero-queue} latencies, the measured
$h_{L2}=\val{amat_rand_hl2_last}$, $h_{L3}=\val{amat_rand_hl3_last}$, $m_D=\val{amat_rand_md_last}$ and the
fitted $t_0$, it predicts \valqty{amat_rand_serial_ns}{\nano\second} per element
(\valqty{amat_rand_serial_dram_ns}{\nano\second} of it DRAM) against the measured
\valqty{cm_rand_max_ns}{\nano\second}, a \val{amat_rand_serial_ratio}-fold overprediction.

Caveats: the L2/L3 terms are only roughly determined. The largest relative errors (random up to
\val{amat_rand_maxrelerr}\%, sequential \val{amat_seq_maxrelerr}\%) are at $N=2^{8}$
(\val{amat_rand_256_relerr}\%, loop overhead) and $N=2^{20}$ (\val{amat_rand_1048576_relerr}\%), where the
footprint equals the L3: only \val{cm_rand_1m_llc_miss} LLC misses but already \val{cm_rand_1m_dtlb_miss}
page walks per element, and the model has no term for walks without DRAM misses
(\qty{2}{\mebi\byte} pages: \valqty{tlbs_huge_1m_ns}{\nano\second}, \cref{tab:tlb-saxpy}).
